\documentclass[11pt,twoside]{book}

% ============================================================================
% PACKAGES
% ============================================================================

% Encodage et langue
\usepackage[utf8]{inputenc}
\usepackage[T1]{fontenc}
\usepackage[english]{babel}

% Mathématiques
\usepackage{amsmath,amssymb,amsthm,amsfonts}
\usepackage{mathtools}
\usepackage{bm}
% \usepackage{dsfont} % Not available

% Mise en page
\usepackage[a4paper,margin=2.5cm,inner=3cm,outer=2.5cm]{geometry}
\usepackage{setspace}
\onehalfspacing
\usepackage{fancyhdr}
\usepackage{titlesec}

% Figures et tableaux
\usepackage{graphicx}
\usepackage{tikz}
\usepackage{tikz-cd}
\usepackage{pgfplots}
\pgfplotsset{compat=1.18}
\usetikzlibrary{arrows.meta,positioning,calc,patterns,decorations.pathmorphing}
\usepackage{booktabs}
\usepackage{multirow}
\usepackage{array}
\usepackage{longtable}

% Références et liens
\usepackage{hyperref}
\hypersetup{
    colorlinks=true,
    linkcolor=blue!70!black,
    citecolor=green!50!black,
    urlcolor=blue!70!black,
    pdftitle={Statistical Estimation with Twin Differential Calculus},
    pdfauthor={},
}
\usepackage{cleveref}

% Algorithmes
\usepackage{algorithm}
\usepackage{algpseudocode}

% Index et glossaire
\usepackage{makeidx}
\makeindex

% Couleurs
\usepackage{xcolor}

% Énumérations
\usepackage{enumitem}

% ============================================================================
% ENVIRONNEMENTS DE THÉORÈMES
% ============================================================================

\theoremstyle{plain}
\newtheorem{theorem}{Theorem}[chapter]
\newtheorem{proposition}[theorem]{Proposition}
\newtheorem{lemma}[theorem]{Lemma}
\newtheorem{corollary}[theorem]{Corollary}

\theoremstyle{definition}
\newtheorem{definition}[theorem]{Definition}
\newtheorem{example}[theorem]{Example}
\newtheorem{algorithm_def}[theorem]{Algorithm}
\newtheorem{principle}[theorem]{Principle}

\theoremstyle{remark}
\newtheorem{remark}[theorem]{Remark}
\newtheorem{note}[theorem]{Note}
\newtheorem{openproblem}{Open Problem}[chapter]

% ============================================================================
% COMMANDES PERSONNALISÉES
% ============================================================================

% Ensembles
\newcommand{\R}{\mathbb{R}}
\newcommand{\N}{\mathbb{N}}
\newcommand{\Z}{\mathbb{Z}}
\newcommand{\C}{\mathbb{C}}
\newcommand{\K}{\mathbb{K}}
\newcommand{\E}{\mathbb{E}}
\newcommand{\Prob}{\mathbb{P}}

% Opérations jumelles
\newcommand{\tplus}{\oplus_g}
\newcommand{\tminus}{\ominus_g}
\newcommand{\ttimes}{\otimes_g}
\newcommand{\tdiv}{\oslash_g}
\newcommand{\Kg}{K_g}
\newcommand{\Kln}{K_{\ln}}
\newcommand{\Klogit}{K_{\mathrm{logit}}}

% Opérateurs
\DeclareMathOperator*{\argmin}{arg\,min}
\DeclareMathOperator*{\argmax}{arg\,max}
\DeclareMathOperator{\Var}{Var}
\DeclareMathOperator{\Cov}{Cov}
\DeclareMathOperator{\diag}{diag}
\DeclareMathOperator{\tr}{tr}
\DeclareMathOperator{\KL}{KL}
\DeclareMathOperator{\Fisher}{I}
\DeclareMathOperator{\sgn}{sgn}
\DeclareMathOperator{\softmax}{softmax}
\DeclareMathOperator{\logit}{logit}

% Calcul différentiel
\newcommand{\dd}{\mathrm{d}}
\newcommand{\pp}[2]{\frac{\partial #1}{\partial #2}}
\newcommand{\ddt}[1]{\frac{\dd #1}{\dd t}}
\newcommand{\Dg}{D_g}
\newcommand{\nablg}{\nabla_g}
\newcommand{\nabnat}{\nabla_{\mathrm{nat}}}

% Espaces
\newcommand{\simplex}{\Delta}
\newcommand{\Rpos}{\R_{>0}}

% ============================================================================
% EN-TÊTES ET PIEDS DE PAGE
% ============================================================================

\pagestyle{fancy}
\fancyhf{}
\fancyhead[LE]{\leftmark}
\fancyhead[RO]{\rightmark}
\fancyfoot[C]{\thepage}
\renewcommand{\headrulewidth}{0.4pt}

% ============================================================================
% TITRE
% ============================================================================

\title{
    \vspace{-2cm}
    {\Huge \textbf{Statistical Estimation with\\[0.3cm] Twin Differential Calculus}}\\[1.5cm]
    {\Large A Unified Geometric Framework for\\[0.2cm] Maximum Likelihood, Bayesian Inference,\\[0.2cm] and Information Geometry}\\[2cm]
    {\large \textit{Comprehensive Book Plan}}
}

\author{
    \textbf{[Author Name]}\\[0.3cm]
    \textit{[Institution]}\\[0.5cm]
    Based on the foundational work:\\
    \textsc{Introduction to Twin Differential Calculus}\\
    by Jocelyn Nembe
}

\date{\today\\[1cm]\small Draft Version 1.0}

% ============================================================================
% DOCUMENT
% ============================================================================

\begin{document}

\frontmatter

\maketitle

% ============================================================================
% PRÉFACE
% ============================================================================

\chapter*{Preface}
\addcontentsline{toc}{chapter}{Preface}

Statistical estimation has long relied on a patchwork of techniques: maximum likelihood estimation operates in parameter space, natural gradient descent exploits Fisher information geometry, and Bayesian methods navigate probability simplices. While each approach has proven powerful, their underlying unity has remained obscured by incompatible formalisms and coordinate-dependent presentations.

This book proposes a radical unification: \textbf{twin differential calculus} provides a single geometric framework that reveals these disparate methods as manifestations of a common principle—\textit{choose your arithmetic according to the problem}.

The key insight is structural: a \textbf{twin space} $\Kg = (K, \tplus, \ttimes)$ is constructed by transporting standard arithmetic through a diffeomorphism $\varphi_g : K \to \R^n$. While the base space may appear globally non-linear, its tangent spaces remain classical vector spaces. This two-level architecture—deformed base, linear tangents—preserves the essence of calculus while adapting to the intrinsic geometry of statistical problems.

\paragraph{Central Thesis.} We demonstrate that:
\begin{itemize}
    \item \textbf{Natural gradient} = twin gradient in logarithmic space $\Kln$
    \item \textbf{Fisher information} = pullback metric induced by $\varphi_g$
    \item \textbf{Exponential families} = twin spaces with $\varphi_g = \nabla A$ (log-partition gradient)
    \item \textbf{EM algorithm} = alternating twin gradient descent
    \item \textbf{MCMC methods} = Hamiltonian dynamics in twin phase space
\end{itemize}

\paragraph{Structure.} The book is organized in six parts:
\begin{enumerate}
    \item \textbf{Foundations} (Chapters 1--3): Motivation, twin calculus review, and information geometry unification
    \item \textbf{Maximum Likelihood Estimation} (Chapters 4--7): MLE reformulation, optimization algorithms, asymptotics, exponential families
    \item \textbf{Bayesian Inference} (Chapters 8--10): Prior/posterior in twin coordinates, MCMC, point estimation
    \item \textbf{Robust and Nonparametric Methods} (Chapters 11--13): M-estimators, kernel methods, hypothesis testing
    \item \textbf{Advanced Models} (Chapters 14--17): Latent variables, time series, deep learning, applications
    \item \textbf{Computation and Perspectives} (Chapters 18--20): Implementation, open problems, conclusion
\end{enumerate}

\paragraph{Prerequisites.} We assume familiarity with:
\begin{itemize}
    \item Graduate-level probability and statistics
    \item Basic differential geometry (tangent spaces, metrics)
    \item Optimization theory (convexity, gradient methods)
    \item Linear algebra and real analysis
\end{itemize}

\paragraph{Acknowledgments.} This work builds directly on Jocelyn Nembe's foundational treatise \textit{Introduction to Twin Differential Calculus}, which established the rigorous mathematical framework underlying our statistical applications.

\vspace{1cm}
\hfill \textit{[Location], \today}

% ============================================================================
% TABLE DES MATIÈRES
% ============================================================================

\tableofcontents

\listoffigures
\addcontentsline{toc}{chapter}{List of Figures}

\listoftables
\addcontentsline{toc}{chapter}{List of Tables}

% ============================================================================
% NOTATION
% ============================================================================

\chapter*{Notation and Conventions}
\addcontentsline{toc}{chapter}{Notation and Conventions}

\section*{Sets and Spaces}

\begin{tabular}{ll}
    $\R, \R^n$ & Real numbers, $n$-dimensional Euclidean space \\
    $\Rpos$ & Positive real numbers $(0, \infty)$ \\
    $\simplex^{d-1}$ & Probability simplex $\{p \in \Rpos^d : \sum_i p_i = 1\}$ \\
    $\Kg = (K, \tplus, \ttimes)$ & Twin space with gauge $\varphi_g$ \\
    $\Kln$ & Logarithmic twin space $(\Rpos, \cdot, \exp)$ \\
    $\Klogit$ & Logit twin space on the simplex \\
    $T_x \Kg$ & Tangent space at $x \in \Kg$ (isomorphic to $\R^n$) \\
    $\mathcal{P}(\mathcal{X})$ & Probability distributions on $\mathcal{X}$ \\
    $\mathcal{M}$ & Statistical manifold \\
\end{tabular}

\section*{Twin Operations}

\begin{tabular}{ll}
    $\varphi_g : K \to \R^n$ & Gauge map (twin isomorphism) \\
    $x \tplus y$ & Twin addition: $\varphi_g^{-1}(\varphi_g(x) + \varphi_g(y))$ \\
    $x \tminus y$ & Twin subtraction: $\varphi_g^{-1}(\varphi_g(x) - \varphi_g(y))$ \\
    $\alpha \ttimes x$ & Twin scalar multiplication: $\varphi_g^{-1}(\alpha \cdot \varphi_g(x))$ \\
    $0_g, 1_g$ & Twin zero and unit: $\varphi_g^{-1}(0), \varphi_g^{-1}(1)$ \\
\end{tabular}

\section*{Differential Calculus}

\begin{tabular}{ll}
    $\Dg f(x)$ & Twin differential of $f$ at $x$ \\
    $\nablg f(x)$ & Twin gradient of $f$ at $x$ \\
    $\nabnat f$ & Natural gradient (Amari) \\
    $D^2_g f(x)$ & Twin Hessian \\
    $\tilde{f} = \varphi_h \circ f \circ \varphi_g^{-1}$ & Conjugated function \\
\end{tabular}

\section*{Statistical Quantities}

\begin{tabular}{ll}
    $p_\theta(x)$ & Probability density/mass function \\
    $\ell(\theta) = \log p_\theta(x)$ & Log-likelihood \\
    $L(\theta) = \prod_i p_\theta(x_i)$ & Likelihood function \\
    $\Fisher(\theta)$ & Fisher information matrix \\
    $\KL(p \| q)$ & Kullback-Leibler divergence \\
    $\theta, \eta$ & Natural and expectation parameters \\
    $A(\theta)$ & Log-partition function \\
    $T(x)$ & Sufficient statistic \\
    $\hat{\theta}_n$ & Maximum likelihood estimator \\
\end{tabular}

\section*{Conventions}

\begin{itemize}
    \item Einstein summation convention: repeated indices are summed
    \item Gradient $\nabla f$ is a column vector; Jacobian $Df$ is a matrix
    \item Log-likelihoods are maximized; losses are minimized
    \item $\stackrel{d}{\to}$ denotes convergence in distribution
    \item $\stackrel{p}{\to}$ denotes convergence in probability
    \item $O_p(\cdot), o_p(\cdot)$ denote stochastic order notation
\end{itemize}

% ============================================================================
% CORPS DU LIVRE
% ============================================================================

\mainmatter

% ############################################################################
% PART I: FOUNDATIONS
% ############################################################################

\part{Foundations}
\label{part:foundations}

% ============================================================================
% CHAPTER 1: MOTIVATIONS
% ============================================================================

\chapter{Motivations: Why Twin Calculus for Statistics?}
\label{ch:motivations}

\begin{quote}
    \textit{``The choice of coordinate system should not be arbitrary—it should reflect the intrinsic geometry of the problem.''}
    
    \hfill --- Shun-ichi Amari
\end{quote}

\section{The Geometry of Statistical Models}
\label{sec:geometry_stat_models}

Statistical models form manifolds: the family $\{p_\theta : \theta \in \Theta\}$ traces out a curved surface in the space of probability distributions. Classical estimation theory often ignores this geometry, treating parameters as Euclidean vectors and applying standard optimization.

\begin{example}[The Simplex Problem]
    Consider estimating a categorical distribution $p = (p_1, \ldots, p_d)$ on the simplex $\simplex^{d-1}$. Standard gradient descent on the log-likelihood
    \[
    \ell(p) = \sum_{i=1}^d n_i \log p_i
    \]
    produces updates $p^{(t+1)} = p^{(t)} + \alpha \nabla \ell(p^{(t)})$ that:
    \begin{enumerate}
        \item May exit the simplex (negative probabilities)
        \item Converge slowly near boundaries
        \item Ignore the multiplicative nature of probability updates
    \end{enumerate}
\end{example}

\section{Limitations of Classical Approaches}
\label{sec:limitations}

\subsection{Coordinate Dependence}

The Euclidean gradient $\nabla f$ depends on the parameterization. For $f : \Rpos \to \R$, the gradient in $\theta$ differs from that in $\eta = \log \theta$:
\[
\nabla_\theta f = f'(\theta), \quad \nabla_\eta f = \theta \cdot f'(\theta).
\]
Which is ``correct''? Neither—both are coordinate artifacts.

\subsection{Constraint Violations}

Parameters often live in constrained domains:
\begin{itemize}
    \item Variances: $\sigma^2 > 0$
    \item Probabilities: $p \in (0,1)$ or $p \in \simplex^{d-1}$
    \item Correlation matrices: positive definite
    \item Covariance matrices: symmetric positive definite
\end{itemize}

Standard methods require projections, barrier functions, or reparameterizations—all coordinate-dependent hacks.

\subsection{Poor Conditioning}

The Hessian $\nabla^2 \ell(\theta)$ can be highly ill-conditioned, especially for:
\begin{itemize}
    \item Hierarchical models (scale separation)
    \item Deep networks (vanishing/exploding gradients)
    \item Near-singular covariance estimation
\end{itemize}

\section{The Twin Calculus Solution}
\label{sec:twin_solution}

Twin differential calculus resolves these issues by a simple principle:

\begin{center}
    \fbox{\parbox{0.8\textwidth}{\centering
        \textbf{Choose the arithmetic to match the problem.}\\[0.2cm]
        Work in a twin space $\Kg$ where constraints become automatic\\
        and the optimization landscape is well-conditioned.
    }}
\end{center}

\subsection{Key Insight: Two-Level Structure}

The resolution rests on a fundamental separation:
\begin{enumerate}
    \item The \textbf{base space} $\Kg = (K, \tplus, \ttimes)$ may have non-standard operations
    \item The \textbf{tangent spaces} $T_x \Kg \cong \R^n$ are always classical vector spaces
    \item The \textbf{differential} $\Dg f(x) : T_x \Kg \to T_{f(x)} K_h$ is always a linear map
\end{enumerate}

This ensures that while parameters live in ``exotic'' spaces, infinitesimal analysis remains classical.

\section{Preview: Unification Results}
\label{sec:preview}

The twin framework reveals deep connections:

\begin{theorem}[Natural Gradient as Twin Gradient]
    \label{thm:preview_natural}
    For the logarithmic twin space $\Kln = (\Rpos, \cdot, \exp)$, and a Fisher
    information of the diagonal form $\Fisher(x) = \diag(1/x_i)$,
    \[
    \nabla_{\ln} f(x) = \nabnat f(x) = x \cdot \nabla f(x),
    \]
    where $\cdot$ is elementwise multiplication. Natural gradient descent is
    gradient descent in $\Kln$.
\end{theorem}

\begin{theorem}[Fisher Metric as Pullback]
    \label{thm:preview_fisher}
    The Fisher information metric $g_{ij}(\theta) = \E\left[\pp{\log p_\theta}{{\theta^i}} \pp{\log p_\theta}{{\theta^j}}\right]$ is the pullback of the Euclidean metric under the score map:
    \[
    g_\theta(u, v) = \langle D\varphi_g(\theta) u, D\varphi_g(\theta) v \rangle.
    \]
\end{theorem}

\begin{theorem}[Exponential Families as Twin Spaces]
    \label{thm:preview_expfam}
    An exponential family $p_\theta(x) = h(x) \exp(\langle \theta, T(x) \rangle - A(\theta))$ has a canonical twin structure with:
    \[
    \varphi_g = \nabla A : \Theta \to \mathcal{H} \quad \text{(natural to expectation parameters)}.
    \]
\end{theorem}

\section{Roadmap}
\label{sec:roadmap}

\begin{figure}[htbp]
    \centering
    \begin{tikzpicture}[
        node distance=1.5cm,
        every node/.style={draw, rounded corners, minimum width=3cm, minimum height=0.8cm, align=center},
        arrow/.style={->, thick}
    ]
        % Part I
        \node (ch1) {Ch.1: Motivations};
        \node (ch2) [below of=ch1] {Ch.2: Twin Spaces};
        \node (ch3) [below of=ch2] {Ch.3: Info Geometry};
        
        % Part II
        \node (ch4) [right of=ch1, xshift=4cm] {Ch.4: MLE};
        \node (ch5) [below of=ch4] {Ch.5: Optimization};
        \node (ch6) [below of=ch5] {Ch.6: Asymptotics};
        \node (ch7) [below of=ch6] {Ch.7: Exp. Families};
        
        % Arrows
        \draw[arrow] (ch1) -- (ch2);
        \draw[arrow] (ch2) -- (ch3);
        \draw[arrow] (ch3) -- (ch4);
        \draw[arrow] (ch4) -- (ch5);
        \draw[arrow] (ch5) -- (ch6);
        \draw[arrow] (ch6) -- (ch7);
        
        % Labels
        \node[draw=none, above of=ch1, yshift=-0.5cm] {\textbf{Part I}};
        \node[draw=none, above of=ch4, yshift=-0.5cm] {\textbf{Part II}};
    \end{tikzpicture}
    \caption{Dependency structure of the first two parts}
    \label{fig:roadmap}
\end{figure}

\section{Notes and References}
\label{sec:ch1_notes}

The geometric view of statistics originates with Fisher (1925) and was systematized by Rao (1945) and Amari (1985). The natural gradient was introduced in \cite{amari1998natural}. Twin differential calculus was developed by Nembe (2025).

% ============================================================================
% CHAPTER 2: TWIN SPACES REVIEW
% ============================================================================

\chapter{Twin Differential Calculus: A Concise Review}
\label{ch:twin_review}

This chapter provides a self-contained summary of twin differential calculus, following the foundational treatment in \cite{nembe2025}. We establish definitions, state the main theorems, and develop the tools needed for statistical applications.

\section{Twin Spaces: Definition and Examples}
\label{sec:twin_definition}

\begin{definition}[Twin Space]
    \label{def:twin_space}
    Let $K \subseteq \R^n$ be an open set and $\varphi_g : K \to \R^n$ a $C^1$-diffeomorphism (the \textbf{gauge map} or \textbf{twin isomorphism}). The \textbf{twin space} generated by $\varphi_g$ is
    \[
    \Kg = (K, \tplus, \ttimes)
    \]
    where the \textbf{twin operations} are defined by transport:
    \begin{align}
        x \tplus y &:= \varphi_g^{-1}(\varphi_g(x) + \varphi_g(y)), \\
        \alpha \ttimes x &:= \varphi_g^{-1}(\alpha \cdot \varphi_g(x)), \quad \alpha \in \R.
    \end{align}
    The \textbf{twin zero} and \textbf{twin unit} are $0_g := \varphi_g^{-1}(0)$ and $1_g := \varphi_g^{-1}(1)$.
\end{definition}

\begin{proposition}[Algebraic Structure]
    $(\Kg, \tplus, \ttimes)$ is a real vector space, isomorphic to $\R^n$ via $\varphi_g$.
\end{proposition}

\subsection{Fundamental Examples}

\begin{example}[Logarithmic Twin Space $\Kln$]
    \label{ex:log_twin}
    Let $K = \Rpos$ and $\varphi_{\ln}(x) = \ln x$. Then:
    \begin{align}
        x \oplus_{\ln} y &= \varphi_{\ln}^{-1}(\ln x + \ln y) = e^{\ln x + \ln y} = x \cdot y, \\
        \alpha \otimes_{\ln} x &= \varphi_{\ln}^{-1}(\alpha \ln x) = e^{\alpha \ln x} = x^\alpha.
    \end{align}
    Thus, $\Kln = (\Rpos, \cdot, \text{power})$: twin addition is multiplication, twin scalar multiplication is exponentiation.
    
    \textbf{Statistical relevance}: Positive parameters (variances, rates, scales).
\end{example}

\begin{example}[Logit Twin Space $\Klogit$]
    \label{ex:logit_twin}
    Let $K = (0,1)$ and $\varphi_{\logit}(p) = \log\frac{p}{1-p}$. Then:
    \begin{align}
        p_1 \oplus_{\logit} p_2 &= \sigma\left(\logit(p_1) + \logit(p_2)\right), \\
        \alpha \otimes_{\logit} p &= \sigma(\alpha \cdot \logit(p)),
    \end{align}
    where $\sigma(x) = 1/(1+e^{-x})$ is the sigmoid function.
    
    \textbf{Statistical relevance}: Binary probabilities, logistic regression.
\end{example}

\begin{example}[Simplex Twin Space]
    \label{ex:simplex_twin}
    For the probability simplex $\simplex^{d-1}$, use $\varphi_{\text{softmax}}^{-1}(u) = \softmax(u)$:
    \[
    \softmax(u)_i = \frac{e^{u_i}}{\sum_j e^{u_j}}.
    \]
    This maps $\R^{d-1}$ (using a reference category) to the interior of the simplex.
    
    \textbf{Statistical relevance}: Categorical distributions, multinomial models.
\end{example}

\begin{example}[Exponential Hierarchy $K_n$]
    \label{ex:exp_hierarchy}
    Define $\varphi_n = \underbrace{\ln \circ \cdots \circ \ln}_{n \text{ times}}$ on appropriate domains:
    \begin{align}
        K_0 &= (\R, +, \cdot) \quad \text{(classical)}, \\
        K_1 &= (\Rpos, \cdot, \exp) = \Kln, \\
        K_2 &= ((e, \infty), \oplus_2, \otimes_2), \\
        &\vdots
    \end{align}
    Each level captures phenomena at different scales of growth.
\end{example}

\section{Twin Differential and the Transfer Theorem}
\label{sec:twin_differential}

\begin{definition}[Twin Differential]
    \label{def:twin_differential}
    Let $f : \Kg \to K_h$ be a function between twin spaces. The \textbf{twin differential} of $f$ at $x \in K$ is the linear map
    \[
    \Dg f(x) : T_x \Kg \to T_{f(x)} K_h
    \]
    defined by
    \[
    \Dg f(x) := D(\varphi_h \circ f \circ \varphi_g^{-1})(\varphi_g(x)),
    \]
    where $D$ denotes the classical (Fréchet) differential.
\end{definition}

\begin{remark}[Two-Level Structure]
    The base spaces $\Kg, K_h$ may have non-standard arithmetic, but:
    \begin{itemize}
        \item Tangent spaces $T_x \Kg \cong \R^n$, $T_{f(x)} K_h \cong \R^m$ are classical
        \item The differential $\Dg f(x)$ is a classical linear map
        \item All calculus rules apply at the tangent level
    \end{itemize}
\end{remark}

\begin{theorem}[Transfer Theorem]
    \label{thm:transfer}
    Let $f : \Kg \to K_h$ be differentiable. Define the \textbf{conjugated function}
    \[
    \tilde{f} := \varphi_h \circ f \circ \varphi_g^{-1} : \R^n \to \R^m.
    \]
    Then:
    \begin{enumerate}
        \item $f$ is twin-differentiable at $x$ if and only if $\tilde{f}$ is differentiable at $\varphi_g(x)$
        \item $\Dg f(x) = D\tilde{f}(\varphi_g(x))$
        \item \textbf{All twin calculus reduces to classical calculus via coordinate transformation}
    \end{enumerate}
    
    \begin{center}
    \begin{tikzcd}
        \Kg \arrow[r, "f"] \arrow[d, "\varphi_g"'] & K_h \arrow[d, "\varphi_h"] \\
        \R^n \arrow[r, "\tilde{f}"'] & \R^m
    \end{tikzcd}
    \end{center}
\end{theorem}

\begin{proof}
    By definition, $\Dg f(x) = D(\varphi_h \circ f \circ \varphi_g^{-1})(\varphi_g(x)) = D\tilde{f}(\varphi_g(x))$. The equivalence of differentiability follows from $\varphi_g, \varphi_h$ being $C^1$-diffeomorphisms.
\end{proof}

\section{Twin Gradient}
\label{sec:twin_gradient}

\begin{definition}[Twin Gradient]
    \label{def:twin_gradient}
    For $F : \Kg \to \R$, the \textbf{twin gradient} is the unique vector $\nablg F(x) \in T_x \Kg \cong \R^n$ satisfying
    \[
    \Dg F(x) \cdot v = \langle \nablg F(x), v \rangle, \quad \forall v \in T_x \Kg,
    \]
    where $\langle \cdot, \cdot \rangle$ is the Euclidean inner product.
\end{definition}

\begin{proposition}[Explicit Formula]
    \label{prop:twin_grad_formula}
    \[
    \nablg F(x) = \nabla(F \circ \varphi_g^{-1})(\varphi_g(x)).
    \]
    The twin gradient equals the classical gradient of the transformed function.
\end{proposition}

\begin{theorem}[Riemannian Interpretation]
    \label{thm:riemannian}
    The twin gradient $\nablg F$ is the Riemannian gradient for the pullback metric
    \[
    g_x(u, v) = \langle D\varphi_g(x) u, D\varphi_g(x) v \rangle.
    \]
    In coordinates, the metric tensor is $G(x) = (D\varphi_g(x))^\top D\varphi_g(x)$.
\end{theorem}

\section{Calculus Rules in Twin Spaces}
\label{sec:calculus_rules}

All classical differentiation rules hold in twin spaces, interpreted at the tangent level.

\begin{theorem}[Chain Rule]
    \label{thm:chain_rule}
    For $f : \Kg \to K_h$ and $g : K_h \to K_k$:
    \[
    D_g(g \circ f)(x) = D_h g(f(x)) \circ \Dg f(x).
    \]
\end{theorem}

\begin{theorem}[Product Rule]
    For $f, g : \Kg \to \R$:
    \[
    \nablg(f \cdot g)(x) = f(x) \nablg g(x) + g(x) \nablg f(x).
    \]
    Note: the product $f \cdot g$ is classical multiplication of real values.
\end{theorem}

\begin{theorem}[Inverse Function Theorem]
    If $\Dg f(x) : \R^n \to \R^n$ is invertible, then $f^{-1}$ exists locally and
    \[
    D_h f^{-1}(f(x)) = (\Dg f(x))^{-1}.
    \]
\end{theorem}

\section{Higher-Order Derivatives}
\label{sec:higher_derivatives}

\begin{definition}[Twin Hessian]
    The second twin derivative $D^2_g f(x)$ is the bilinear form
    \[
    D^2_g f(x) := D^2(\varphi_h \circ f \circ \varphi_g^{-1})(\varphi_g(x)).
    \]
    For $F : \Kg \to \R$, the twin Hessian matrix is
    \[
    H_g F(x) = \nabla^2(F \circ \varphi_g^{-1})(\varphi_g(x)).
    \]
\end{definition}

\begin{remark}
    Higher derivatives involve the curvature of the gauge $\varphi_g$. Expressions may contain terms proportional to $D^2 \varphi_g$, reflecting coordinate effects.
\end{remark}

\section{Summary: The Twin Calculus Toolkit}
\label{sec:toolkit_summary}

\begin{table}[htbp]
    \centering
    \caption{Twin calculus operations and their classical equivalents}
    \label{tab:toolkit}
    \begin{tabular}{lll}
        \toprule
        \textbf{Twin Operation} & \textbf{Definition} & \textbf{Classical Equivalent} \\
        \midrule
        Addition $x \tplus y$ & $\varphi_g^{-1}(\varphi_g(x) + \varphi_g(y))$ & $u + v$ in $\varphi_g$-coords \\
        Scalar mult. $\alpha \ttimes x$ & $\varphi_g^{-1}(\alpha \varphi_g(x))$ & $\alpha u$ in $\varphi_g$-coords \\
        Differential $\Dg f(x)$ & $D\tilde{f}(\varphi_g(x))$ & Classical $D\tilde{f}$ \\
        Gradient $\nablg F(x)$ & $\nabla \tilde{F}(\varphi_g(x))$ & Classical $\nabla \tilde{F}$ \\
        Hessian $H_g F(x)$ & $\nabla^2 \tilde{F}(\varphi_g(x))$ & Classical $\nabla^2 \tilde{F}$ \\
        \bottomrule
    \end{tabular}
\end{table}

\section{Notes and References}
\label{sec:ch2_notes}

The complete theory of twin differential calculus is developed in \cite{nembe2025}. The Transfer Theorem (Theorem \ref{thm:transfer}) is the central result, reducing all twin computations to classical calculus. For the exponential hierarchy, see \cite{grossman1972}.

% ============================================================================
% CHAPTER 3: INFORMATION GEOMETRY
% ============================================================================

\chapter{Unification with Information Geometry}
\label{ch:info_geometry}

This chapter establishes the deep connection between twin calculus and information geometry. We prove that the Fisher-Rao metric arises as a pullback metric, that natural gradient descent is twin gradient descent in $\Kln$, and that exponential families have canonical twin structure.

\section{The Fisher Information Metric}
\label{sec:fisher_metric}

\begin{definition}[Fisher Information]
    For a statistical model $\{p_\theta : \theta \in \Theta\}$, the \textbf{Fisher information matrix} is
    \[
    \Fisher(\theta)_{ij} = \E_{p_\theta}\left[\pp{\log p_\theta}{\theta^i} \pp{\log p_\theta}{\theta^j}\right] = -\E_{p_\theta}\left[\frac{\partial^2 \log p_\theta}{\partial \theta^i \partial \theta^j}\right].
    \]
\end{definition}

\begin{definition}[Fisher-Rao Metric]
    The \textbf{Fisher-Rao metric} on the statistical manifold $\mathcal{M} = \{p_\theta\}$ is
    \[
    g_\theta^{FR}(u, v) = u^\top \Fisher(\theta) v = \sum_{i,j} \Fisher(\theta)_{ij} u^i v^j.
    \]
\end{definition}

\begin{theorem}[Fisher Metric as Pullback]
    \label{thm:fisher_pullback}
    Let $\varphi_g : \Theta \to \R^d$ be a diffeomorphism. Define the pullback metric
    \[
    g_\theta(u, v) := \langle D\varphi_g(\theta) u, D\varphi_g(\theta) v \rangle.
    \]
    Take for $\varphi_g$ the log-likelihood embedding
    $\iota : \theta \mapsto \log p_\theta(\cdot) \in L^2(p_\theta)$, whose
    differential is the score, $D\iota(\theta) e_i = s_i(\theta) :=
    \partial_{\theta^i} \log p_\theta$. Then the pullback of the $L^2(p_\theta)$
    inner product coincides with the Fisher--Rao metric:
    \[
    g_\theta = g_\theta^{FR}.
    \]
\end{theorem}

\begin{proof}
    The differential of the embedding sends the coordinate vector $e_i$ to the
    $i$-th score component, $D\iota(\theta) e_i = s_i(\theta)$. Hence
    \begin{align}
        g_\theta(e_i, e_j) &= \langle D\iota(\theta) e_i, D\iota(\theta) e_j \rangle_{L^2(p_\theta)} \\
        &= \E_{p_\theta}[s_i(\theta) s_j(\theta)] \\
        &= \Fisher(\theta)_{ij}, \qedhere
    \end{align}
    using $\E_{p_\theta}[s_i s_j] = \Fisher(\theta)_{ij}$ by definition of the
    Fisher information.
\end{proof}

\section{Natural Gradient as Twin Gradient}
\label{sec:natural_gradient}

\begin{definition}[Natural Gradient (Amari)]
    The \textbf{natural gradient} of $F : \Theta \to \R$ is
    \[
    \nabnat F(\theta) = \Fisher(\theta)^{-1} \nabla F(\theta).
    \]
    This is the Riemannian gradient for the Fisher-Rao metric.
\end{definition}

\begin{theorem}[Natural Gradient = Twin Gradient in $\Kln$]
    \label{thm:natural_twin}
    For the logarithmic twin space $\Kln = (\Rpos, \cdot, \exp)$ with $\varphi_{\ln}(x) = \ln x$:
    \[
    \nabla_{\ln} F(x) = \nabla(F \circ \exp)(\ln x) = x \cdot \nabla F(x) = \nabnat F(x)
    \]
    whenever the Fisher information has the diagonal form $\Fisher = \diag(1/x_i)$,
    so that $\Fisher^{-1} = \diag(x_i)$ and $\nabnat F = \Fisher^{-1}\nabla F = x \cdot \nabla F$.
    
    More generally, on the probability simplex $\simplex^{d-1}$ with KL geometry:
    \[
    \nabla_{\ln} F(p) = (p_1 \partial_1 F(p), \ldots, p_d \partial_d F(p)) = \nabnat F(p).
    \]
\end{theorem}

\begin{proof}
    Let $\tilde{F}(u) = F(e^u)$ where $u = \ln p$ (componentwise). Then:
    \begin{align}
        \nabla_{\ln} F(p) &= \nabla \tilde{F}(\ln p) = \nabla_u F(e^u)\big|_{u=\ln p} \\
        &= \left(\pp{F}{u_i}\right)_{i=1}^d = \left(p_i \pp{F}{p_i}\right)_{i=1}^d.
    \end{align}
    For the simplex with Fisher metric $\Fisher(p)_{ij} = \delta_{ij}/p_i$:
    \[
    \nabnat F(p) = \Fisher(p)^{-1} \nabla F(p) = (p_1 \partial_1 F, \ldots, p_d \partial_d F). \qedhere
    \]
\end{proof}

\begin{remark}[Scope of the identity]
    The equality $\nabla_{\ln} F = \nabnat F$ is not automatic: it requires the
    Fisher information to be the diagonal multiplicative form $\diag(1/x_i)$
    (as for the simplex/categorical and Poisson-count models above), for then
    $\Fisher^{-1}\nabla F = x \cdot \nabla F$. The twin gradient $\nabla_{\ln} F$
    is always the Euclidean gradient of $F$ in the logarithmic chart
    $u = \ln x$; it coincides with Amari's natural gradient precisely when the
    pullback metric $G = (D\varphi_{\ln})^\top D\varphi_{\ln} = \diag(1/x_i^2)$
    of the chart agrees, up to the chart's own Jacobian, with the statistical
    Fisher metric. For a general model these two metrics differ and twin-log
    descent is a preconditioned---but not the natural-gradient---method.
\end{remark}

\begin{corollary}[Multiplicative Weight Updates]
    Twin gradient descent on the simplex:
    \[
    p^{(t+1)} = p^{(t)} \otimes_{\ln} (\alpha \nabla_{\ln} F(p^{(t)}))^{-1}
    \]
    corresponds to the multiplicative update
    \[
    p_i^{(t+1)} \propto p_i^{(t)} \cdot \exp(-\alpha \partial_i F(p^{(t)})).
    \]
    This is the exponentiated gradient / mirror descent algorithm.
\end{corollary}

\section{Exponential Families as Twin Spaces}
\label{sec:exp_families_twin}

\begin{definition}[Exponential Family]
    A family $\{p_\theta\}$ is an \textbf{exponential family} if
    \[
    p_\theta(x) = h(x) \exp(\langle \theta, T(x) \rangle - A(\theta)),
    \]
    where $\theta \in \Theta \subseteq \R^d$ (natural parameters), $T : \mathcal{X} \to \R^d$ (sufficient statistics), and $A(\theta) = \log \int h(x) e^{\langle \theta, T(x) \rangle} dx$ (log-partition).
\end{definition}

\begin{theorem}[Canonical Twin Structure]
    \label{thm:exp_family_twin}
    An exponential family has a canonical twin space structure with:
    \begin{enumerate}
        \item \textbf{Gauge map}: $\varphi_g = \nabla A : \Theta \to \mathcal{H}$
        \item \textbf{Expectation parameters}: $\eta = \nabla A(\theta) = \E_{p_\theta}[T(X)]$
        \item \textbf{Fisher information}: $\Fisher(\theta) = \nabla^2 A(\theta) = \Cov_{p_\theta}(T(X))$
        \item \textbf{Duality}: $\theta \leftrightarrow \eta$ are Legendre dual coordinates
    \end{enumerate}
\end{theorem}

\begin{proof}
    The log-partition function $A(\theta)$ is convex, so $\nabla A$ is a diffeomorphism onto its image. The Hessian relation follows from:
    \[
    \pp{^2 A}{\theta^i \partial \theta^j} = \Cov_{p_\theta}(T_i(X), T_j(X)) = \Fisher(\theta)_{ij}.
    \]
    Legendre duality: defining $A^*(\eta) = \sup_\theta \{\langle \theta, \eta \rangle - A(\theta)\}$, we have $\theta = \nabla A^*(\eta)$.
\end{proof}

\begin{example}[Gaussian Family]
    For $p_\theta(x) = \frac{1}{\sqrt{2\pi\sigma^2}} \exp\left(-\frac{(x-\mu)^2}{2\sigma^2}\right)$:
    \begin{itemize}
        \item Natural parameters: $\theta = (\mu/\sigma^2, -1/(2\sigma^2))$
        \item Sufficient statistics: $T(x) = (x, x^2)$
        \item Expectation parameters: $\eta = (\mu, \mu^2 + \sigma^2)$
        \item Twin space: $\Kg$ with $\varphi_g(\theta) = \eta$
    \end{itemize}
\end{example}

\section{The Dually Flat Structure}
\label{sec:dual_flat}

\begin{definition}[Dual Connections]
    On an exponential family manifold, there exist two affine connections:
    \begin{itemize}
        \item \textbf{$e$-connection} ($\nabla^{(e)}$): flat in $\theta$-coordinates
        \item \textbf{$m$-connection} ($\nabla^{(m)}$): flat in $\eta$-coordinates
    \end{itemize}
    These are dual with respect to the Fisher metric.
\end{definition}

\begin{theorem}[Pythagorean Theorem for KL Divergence]
    \label{thm:pythagorean}
    For $p, q, r$ in an exponential family, if the $m$-geodesic from $p$ to $q$ is orthogonal to the $e$-geodesic from $q$ to $r$:
    \[
    \KL(p \| r) = \KL(p \| q) + \KL(q \| r).
    \]
\end{theorem}

\begin{corollary}[Projection Theorem]
    The MLE is the $m$-projection of the empirical distribution onto the model manifold:
    \[
    \hat{\theta}_{MLE} = \argmin_{\theta \in \Theta} \KL(\hat{p}_n \| p_\theta).
    \]
\end{corollary}

\section{Information Geometry in Twin Coordinates}
\label{sec:info_geom_twin}

\begin{table}[htbp]
    \centering
    \caption{Correspondence between information geometry and twin calculus}
    \label{tab:info_twin}
    \begin{tabular}{ll}
        \toprule
        \textbf{Information Geometry} & \textbf{Twin Calculus} \\
        \midrule
        Fisher metric $g^{FR}$ & Pullback metric $g_x(u,v) = \langle D\varphi_g u, D\varphi_g v \rangle$ \\
        Natural gradient $\nabnat F$ & Twin gradient $\nablg[\ln] F$ \\
        $e$-connection (flat in $\theta$) & Classical calculus in $\theta$-space \\
        $m$-connection (flat in $\eta$) & Classical calculus in $\eta$-space (twin) \\
        Dual coordinates $(\theta, \eta)$ & Gauge map $\varphi_g = \nabla A$ \\
        KL divergence & Bregman divergence in twin coords \\
        \bottomrule
    \end{tabular}
\end{table}

\section{Notes and References}
\label{sec:ch3_notes}

Information geometry was founded by Rao (1945), Chentsov (1972), and systematized by Amari (1985, 2000). The natural gradient for neural networks appears in \cite{amari1998natural}. The connection to twin calculus is developed in \cite{nembe2025}.

% ############################################################################
% PART II: MAXIMUM LIKELIHOOD ESTIMATION
% ############################################################################

\part{Maximum Likelihood Estimation}
\label{part:mle}

% ============================================================================
% CHAPTER 4: MLE REFORMULATION
% ============================================================================

\chapter{Maximum Likelihood in Twin Spaces}
\label{ch:mle}

\section{Classical MLE: A Brief Review}
\label{sec:classical_mle}

Given i.i.d.\ observations $X_1, \ldots, X_n \sim p_\theta$, the \textbf{maximum likelihood estimator} is
\[
\hat{\theta}_n = \argmax_{\theta \in \Theta} L(\theta) = \argmax_{\theta \in \Theta} \prod_{i=1}^n p_\theta(X_i) = \argmax_{\theta \in \Theta} \sum_{i=1}^n \log p_\theta(X_i).
\]

\section{MLE in the Logarithmic Twin Space}
\label{sec:mle_kln}

For positive parameters $\theta \in \Rpos^d$, work in $\Kln$:

\begin{definition}[Twin Log-Likelihood]
    The \textbf{twin log-likelihood} in $\Kln$ is
    \[
    \tilde{\ell}(u) := \ell(e^u) = \sum_{i=1}^n \log p_{e^u}(X_i), \quad u = \ln \theta.
    \]
\end{definition}

\begin{proposition}[Twin Score Equation]
    The MLE satisfies the twin score equation:
    \[
    \nabla_{\ln} \ell(\theta) = \theta \odot \nabla \ell(\theta) = 0,
    \]
    where $\odot$ denotes elementwise multiplication. Equivalently, in $u$-coordinates:
    \[
    \nabla \tilde{\ell}(u) = 0.
    \]
\end{proposition}

\section{MLE on the Probability Simplex}
\label{sec:mle_simplex}

For categorical data with $p \in \simplex^{d-1}$:

\begin{theorem}[Simplex MLE in Twin Coordinates]
    The log-likelihood $\ell(p) = \sum_{i=1}^d n_i \log p_i$ in softmax coordinates $p = \softmax(u)$ becomes:
    \[
    \tilde{\ell}(u) = \sum_{i=1}^d n_i u_i - n \log\left(\sum_j e^{u_j}\right).
    \]
    This is concave in $u$, with gradient
    \[
    \nabla \tilde{\ell}(u) = n \cdot (\hat{p}^{emp} - p(u)),
    \]
    where $\hat{p}^{emp}_i = n_i/n$ is the empirical distribution.
\end{theorem}

\section{Constrained MLE via Twin Spaces}
\label{sec:constrained_mle}

\begin{principle}[Constraint Elimination]
    Choose $\varphi_g$ such that constraints in $\theta$-space become automatic:
    \begin{itemize}
        \item $\theta > 0$: Use $\varphi_{\ln}$, work in $u \in \R$
        \item $\theta \in (0,1)$: Use $\varphi_{\logit}$, work in $u \in \R$
        \item $\theta \in \simplex$: Use $\varphi_{\softmax}^{-1}$, work in $u \in \R^{d-1}$
        \item $\Sigma \succ 0$: Use Cholesky $\Sigma = LL^\top$, work in $L$
    \end{itemize}
\end{principle}

\section{Notes and References}
\label{sec:ch4_notes}

% ============================================================================
% CHAPTER 5: OPTIMIZATION ALGORITHMS
% ============================================================================

\chapter{Optimization Algorithms in Twin Spaces}
\label{ch:optimization}

\section{Twin Gradient Descent}
\label{sec:twin_gd}

\begin{algorithm_def}[Twin Gradient Descent]
    \label{alg:twin_gd}
    Given objective $F : \Kg \to \R$, step size $\alpha > 0$, initial $x_0 \in K$:
    \begin{algorithmic}[1]
        \For{$t = 0, 1, 2, \ldots$}
            \State $u_t \gets \varphi_g(x_t)$ \Comment{Transform to linear space}
            \State $g_t \gets \nabla(F \circ \varphi_g^{-1})(u_t)$ \Comment{Classical gradient}
            \State $u_{t+1} \gets u_t - \alpha g_t$ \Comment{Euclidean update}
            \State $x_{t+1} \gets \varphi_g^{-1}(u_{t+1})$ \Comment{Map back to $\Kg$}
        \EndFor
    \end{algorithmic}
    
    Equivalently, in twin notation:
    \[
    x_{t+1} = x_t \tminus (\alpha \ttimes \nablg F(x_t)).
    \]
\end{algorithm_def}

\begin{theorem}[Convergence of Twin GD]
    \label{thm:twin_gd_convergence}
    Let $\tilde{F} = F \circ \varphi_g^{-1}$ be $L$-smooth and $\mu$-strongly convex. Then twin gradient descent with step size $\alpha \in (0, 1/L]$ converges linearly:
    \[
    \|\varphi_g(x_t) - \varphi_g(x^*)\| \leq (1 - \alpha\mu)^t \|\varphi_g(x_0) - \varphi_g(x^*)\|.
    \]
    (More generally, for $\alpha \in (0, 2/L)$ the contraction factor is
    $\max\{|1-\alpha\mu|, |1-\alpha L|\} < 1$; it equals $1-\alpha\mu$ on the
    stated range $\alpha \leq 1/L$.)
\end{theorem}

\section{Twin Newton-Raphson}
\label{sec:twin_newton}

\begin{algorithm_def}[Twin Newton Method]
    \[
    u_{t+1} = u_t - [H_g F(x_t)]^{-1} \nablg F(x_t),
    \]
    where $H_g F(x_t) = \nabla^2 \tilde{F}(u_t)$ is the twin Hessian.
\end{algorithm_def}

\begin{theorem}[Quadratic Convergence]
    Near a non-degenerate optimum, twin Newton converges quadratically:
    \[
    \|u_{t+1} - u^*\| \leq C \|u_t - u^*\|^2.
    \]
\end{theorem}

\section{Adaptive Methods: Twin Adam}
\label{sec:twin_adam}

\begin{algorithm_def}[Twin Adam]
    Maintain moment estimates in $u$-space:
    \begin{align}
        m_t &= \beta_1 m_{t-1} + (1-\beta_1) g_t, \\
        v_t &= \beta_2 v_{t-1} + (1-\beta_2) g_t^2, \\
        u_{t+1} &= u_t - \alpha \frac{\hat{m}_t}{\sqrt{\hat{v}_t} + \epsilon},
    \end{align}
    where $\hat{m}_t, \hat{v}_t$ are bias-corrected moments and $g_t = \nabla \tilde{F}(u_t)$.
\end{algorithm_def}

\section{Numerical Examples}
\label{sec:numerical_examples}

\begin{table}[htbp]
    \centering
    \caption{Convergence comparison: Twin GD vs Classical GD}
    \label{tab:convergence}
    \begin{tabular}{lccc}
        \toprule
        \textbf{Problem} & \textbf{Classical GD} & \textbf{Twin GD ($g = \ln$)} & \textbf{Speedup} \\
        \midrule
        Portfolio optimization & 87 iterations & 12 iterations & $7.25\times$ \\
        Logistic regression & 1500 iterations & 200 iterations & $7.5\times$ \\
        Neural net (BatchNorm scales) & -- & 40\% variance reduction & -- \\
        \bottomrule
    \end{tabular}
\end{table}

\section{Notes and References}
\label{sec:ch5_notes}

% ============================================================================
% CHAPTER 6: ASYMPTOTIC THEORY
% ============================================================================

\chapter{Asymptotic Theory in Twin Coordinates}
\label{ch:asymptotics}

\section{Consistency}
\label{sec:consistency}

\begin{theorem}[Twin Consistency]
    Under standard regularity conditions, the MLE $\hat{\theta}_n$ is consistent:
    \[
    \varphi_g(\hat{\theta}_n) \stackrel{p}{\to} \varphi_g(\theta_0).
    \]
\end{theorem}

\section{Asymptotic Normality}
\label{sec:asymp_normality}

\begin{theorem}[Central Limit Theorem in Twin Coordinates]
    \label{thm:twin_clt}
    Under regularity conditions:
    \[
    \sqrt{n}(\varphi_g(\hat{\theta}_n) - \varphi_g(\theta_0)) \stackrel{d}{\to} \mathcal{N}(0, \Fisher_g(\theta_0)^{-1}),
    \]
    where $\Fisher_g(\theta) = (D\varphi_g(\theta))^{-\top} \Fisher(\theta) (D\varphi_g(\theta))^{-1}$ is the Fisher information in $g$-coordinates. Equivalently, by the delta method the asymptotic covariance is $\Fisher_g(\theta_0)^{-1} = D\varphi_g(\theta_0)\, \Fisher(\theta_0)^{-1} (D\varphi_g(\theta_0))^\top$.
\end{theorem}

\section{Cramér-Rao Bound}
\label{sec:cramer_rao}

\begin{theorem}[Twin Cramér-Rao Bound]
    For any unbiased estimator $\hat{\psi}$ of $\psi(\theta) = \varphi_g(\theta)$:
    \[
    \Var(\hat{\psi}) \geq \Fisher_g(\theta)^{-1}.
    \]
    The MLE achieves this bound asymptotically.
\end{theorem}

\section{Confidence Regions}
\label{sec:confidence}

\begin{definition}[Twin Confidence Ellipsoid]
    An asymptotic $(1-\alpha)$-confidence region for $\theta$ is
    \[
    \left\{\theta : n (\varphi_g(\hat{\theta}_n) - \varphi_g(\theta))^\top \hat{\Fisher}_g (\varphi_g(\hat{\theta}_n) - \varphi_g(\theta)) \leq \chi^2_{d,\alpha}\right\}.
    \]
\end{definition}

\section{Notes and References}
\label{sec:ch6_notes}

% ============================================================================
% CHAPTER 7: EXPONENTIAL FAMILIES
% ============================================================================

\chapter{Exponential Families and Generalized Linear Models}
\label{ch:exp_families}

\section{Exponential Family MLE in Twin Coordinates}
\label{sec:exp_mle}

\begin{theorem}[MLE for Exponential Families]
    For an exponential family $p_\theta(x) = h(x) \exp(\langle \theta, T(x) \rangle - A(\theta))$:
    \begin{enumerate}
        \item The MLE satisfies $\nabla A(\hat{\theta}) = \bar{T}_n := \frac{1}{n}\sum_i T(X_i)$
        \item In expectation parameters: $\hat{\eta} = \bar{T}_n$ (moment matching)
        \item The twin gradient is $\nablg \ell(\theta) = n(\bar{T}_n - \nabla A(\theta))$
    \end{enumerate}
\end{theorem}

\section{Generalized Linear Models}
\label{sec:glm}

\begin{definition}[GLM in Twin Form]
    A GLM with canonical link has:
    \begin{itemize}
        \item Linear predictor: $\eta_i = x_i^\top \beta$
        \item Link: $g(\mu_i) = \eta_i$ where $\mu_i = \E[Y_i]$
        \item Response distribution: exponential family
    \end{itemize}
    The canonical link satisfies $\eta = \theta$ (natural parameter), making the model a twin space with $\varphi_g = g^{-1}$.
\end{definition}

\section{Graphical Models}
\label{sec:graphical}

\begin{theorem}[Exponential Family Graphical Models]
    Undirected graphical models with log-linear parameterization:
    \[
    p_\theta(x) \propto \exp\left(\sum_{C \in \mathcal{C}} \theta_C \phi_C(x_C)\right)
    \]
    form exponential families where the twin structure encodes conditional independence.
\end{theorem}

\section{Notes and References}
\label{sec:ch7_notes}

% ############################################################################
% PART III: BAYESIAN INFERENCE
% ############################################################################

\part{Bayesian Inference}
\label{part:bayesian}

% ============================================================================
% CHAPTER 8: BAYESIAN FRAMEWORK
% ============================================================================

\chapter{Bayesian Inference in Twin Spaces}
\label{ch:bayesian}

\section{Prior and Posterior Transformation}
\label{sec:prior_posterior}

\begin{theorem}[Posterior in Twin Coordinates]
    Let $\pi(\theta)$ be a prior on $\Theta$ and $\tilde{\pi}(u) = \pi(\varphi_g^{-1}(u)) |D\varphi_g^{-1}(u)|$ the induced prior in $u$-space. Then the posterior transforms as:
    \[
    \tilde{\pi}(u | X) \propto \tilde{\pi}(u) \cdot L(\varphi_g^{-1}(u); X).
    \]
\end{theorem}

\section{Conjugate Priors as Twin Structures}
\label{sec:conjugate}

\begin{theorem}[Conjugate Families in Exponential Families]
    For an exponential family, the conjugate prior is
    \[
    \pi(\theta | \alpha, \nu) \propto \exp(\langle \theta, \alpha \rangle - \nu A(\theta)),
    \]
    which has the same exponential family form. In twin coordinates $\eta = \nabla A(\theta)$, the posterior mean is a convex combination:
    \[
    \hat{\eta}_{Bayes} = \frac{\nu}{\nu + n} \cdot \frac{\alpha}{\nu} + \frac{n}{\nu + n} \cdot \bar{T}_n.
    \]
\end{theorem}

\section{Jeffreys Prior}
\label{sec:jeffreys}

\begin{definition}[Jeffreys Prior]
    The \textbf{Jeffreys prior} is
    \[
    \pi^J(\theta) \propto \sqrt{\det \Fisher(\theta)}.
    \]
    It is invariant under reparameterization.
\end{definition}

\begin{theorem}[Jeffreys Prior in Twin Coordinates]
    In $u = \varphi_g(\theta)$ coordinates, the Jeffreys prior becomes:
    \[
    \tilde{\pi}^J(u) \propto \sqrt{\det \Fisher_g(u)} = \sqrt{\det(\tilde{\Fisher}(u))},
    \]
    where $\tilde{\Fisher}(u) = \Fisher(\varphi_g^{-1}(u))$ in appropriate coordinates.
\end{theorem}

\section{Notes and References}
\label{sec:ch8_notes}

% ============================================================================
% CHAPTER 9: MCMC METHODS
% ============================================================================

\chapter{MCMC Methods in Twin Coordinates}
\label{ch:mcmc}

\section{Metropolis-Hastings in $\varphi_g$-Coordinates}
\label{sec:mh_twin}

\begin{algorithm_def}[Twin Metropolis-Hastings]
    \begin{enumerate}
        \item Propose $u' \sim q(u' | u_t)$ in $u$-space
        \item Accept with probability $\alpha = \min\left(1, \frac{\tilde{\pi}(u' | X) q(u_t | u')}{\tilde{\pi}(u_t | X) q(u' | u_t)}\right)$
        \item Set $\theta_{t+1} = \varphi_g^{-1}(u_{t+1})$
    \end{enumerate}
\end{algorithm_def}

\section{Hamiltonian Monte Carlo}
\label{sec:hmc}

\begin{theorem}[Twin Hamiltonian Dynamics]
    Define the Hamiltonian in $u$-space:
    \[
    H(u, p) = -\log \tilde{\pi}(u | X) + \frac{1}{2}p^\top M^{-1} p.
    \]
    Hamilton's equations are:
    \begin{align}
        \dot{u} &= M^{-1} p, \\
        \dot{p} &= \nabla \log \tilde{\pi}(u | X).
    \end{align}
    The symplectic structure is preserved under leapfrog integration.
\end{theorem}

\section{Variational Inference}
\label{sec:vi}

\begin{theorem}[ELBO in Twin Coordinates]
    For variational family $q_\phi(u)$ in $u$-space:
    \[
    \mathcal{L}(\phi) = \E_{q_\phi}[\log \tilde{\pi}(u | X)] + H(q_\phi),
    \]
    where $H(q_\phi) = -\E_{q_\phi}[\log q_\phi]$ is the entropy. The natural gradient update is:
    \[
    \phi_{t+1} = \phi_t + \alpha \Fisher_{q}(\phi)^{-1} \nabla_\phi \mathcal{L}(\phi).
    \]
\end{theorem}

\section{Notes and References}
\label{sec:ch9_notes}

% ============================================================================
% CHAPTER 10: POINT ESTIMATION
% ============================================================================

\chapter{Bayesian Point Estimation}
\label{ch:point_estimation}

\section{Bayes Estimators under Twin Loss}
\label{sec:bayes_estimators}

\begin{definition}[Twin Loss Function]
    The \textbf{twin squared error loss} is
    \[
    L_g(\theta, \hat{\theta}) = \|\varphi_g(\theta) - \varphi_g(\hat{\theta})\|^2.
    \]
\end{definition}

\begin{theorem}[Bayes Estimator for Twin Loss]
    The Bayes estimator minimizing posterior expected twin loss is:
    \[
    \hat{\theta}^{Bayes}_g = \varphi_g^{-1}\left(\E[\varphi_g(\theta) | X]\right).
    \]
    This is the posterior mean in $u$-coordinates, mapped back.
\end{theorem}

\section{Credible Regions}
\label{sec:credible}

\begin{definition}[Twin HPD Region]
    The highest posterior density region in $u$-space is
    \[
    C_\alpha = \{u : \tilde{\pi}(u | X) \geq c_\alpha\},
    \]
    mapped back to $\theta$-space via $\varphi_g^{-1}$.
\end{definition}

\section{Notes and References}
\label{sec:ch10_notes}

% ############################################################################
% PART IV: ROBUST AND NONPARAMETRIC METHODS
% ############################################################################

\part{Robust and Nonparametric Methods}
\label{part:robust}

% ============================================================================
% CHAPTER 11: ROBUST ESTIMATION
% ============================================================================

\chapter{Robust Estimation in Twin Spaces}
\label{ch:robust}

\section{M-Estimators in $\Kg$}
\label{sec:m_estimators}

\begin{definition}[Twin M-Estimator]
    A \textbf{twin M-estimator} minimizes
    \[
    \hat{\theta}_\rho = \argmin_{\theta \in K} \sum_{i=1}^n \rho(\varphi_g(X_i), \theta),
    \]
    where $\rho : \R \times K \to \R$ is a loss function.
\end{definition}

\section{Influence Functions}
\label{sec:influence}

\begin{definition}[Twin Influence Function]
    The influence function in $g$-coordinates is
    \[
    IF_g(x; T, F) = \lim_{\epsilon \to 0} \frac{\varphi_g(T((1-\epsilon)F + \epsilon \delta_x)) - \varphi_g(T(F))}{\epsilon}.
    \]
\end{definition}

\section{Outlier Handling}
\label{sec:outliers}

\begin{theorem}[Robustness via Gauge Choice]
    For heavy-tailed data, the logarithmic gauge $\varphi_{\ln}$ provides natural robustness: large observations are compressed, reducing their influence on the estimator.
\end{theorem}

\section{Notes and References}
\label{sec:ch11_notes}

% ============================================================================
% CHAPTER 12: NONPARAMETRIC METHODS
% ============================================================================

\chapter{Nonparametric Estimation in Twin Spaces}
\label{ch:nonparametric}

\section{Kernel Density Estimation}
\label{sec:kde}

\begin{definition}[Twin KDE]
    The kernel density estimator in $u$-space is
    \[
    \tilde{f}_h(u) = \frac{1}{nh^d} \sum_{i=1}^n K\left(\frac{u - \varphi_g(X_i)}{h}\right),
    \]
    mapped back to density on $K$ via $f_h(x) = \tilde{f}_h(\varphi_g(x)) |D\varphi_g(x)|$.
\end{definition}

\section{Twin Splines}
\label{sec:splines}

\section{Functional Estimation}
\label{sec:functional}

\section{Notes and References}
\label{sec:ch12_notes}

% ============================================================================
% CHAPTER 13: HYPOTHESIS TESTING
% ============================================================================

\chapter{Hypothesis Testing in Twin Coordinates}
\label{ch:testing}

\section{Likelihood Ratio Tests}
\label{sec:lrt}

\begin{theorem}[Wilks' Theorem in Twin Coordinates]
    Under $H_0 : \theta \in \Theta_0$:
    \[
    -2\log \Lambda_n = -2\log \frac{\sup_{\theta \in \Theta_0} L(\theta)}{\sup_{\theta \in \Theta} L(\theta)} \stackrel{d}{\to} \chi^2_r,
    \]
    where $r = \dim(\Theta) - \dim(\Theta_0)$. This is invariant under twin reparameterization.
\end{theorem}

\section{Wald and Score Tests}
\label{sec:wald_score}

\begin{theorem}[Twin Wald Test]
    The Wald statistic in $g$-coordinates:
    \[
    W_g = n(\varphi_g(\hat{\theta}) - \varphi_g(\theta_0))^\top \hat{\Fisher}_g (\varphi_g(\hat{\theta}) - \varphi_g(\theta_0)) \stackrel{d}{\to} \chi^2_d.
    \]
\end{theorem}

\section{Notes and References}
\label{sec:ch13_notes}

% ############################################################################
% PART V: ADVANCED MODELS
% ############################################################################

\part{Advanced Models}
\label{part:advanced}

% ============================================================================
% CHAPTER 14: LATENT VARIABLES
% ============================================================================

\chapter{Latent Variable Models and the EM Algorithm}
\label{ch:em}

\section{EM Algorithm in Twin Spaces}
\label{sec:em_twin}

\begin{theorem}[Twin EM Algorithm]
    \label{thm:twin_em}
    For latent variable models $p_\theta(x) = \int p_\theta(x, z) dz$:
    
    \textbf{E-step}: Compute $Q(\theta | \theta^{(t)}) = \E_{Z|X,\theta^{(t)}}[\log p_\theta(X, Z)]$
    
    \textbf{M-step in twin space}: Update via twin gradient ascent on $Q$:
    \[
    \theta^{(t+1)} = \theta^{(t)} \tplus (\alpha \ttimes \nablg Q(\theta^{(t)} | \theta^{(t)})).
    \]
    
    For exponential families, this reduces to moment matching in twin coordinates.
\end{theorem}

\section{Mixture Models}
\label{sec:mixtures}

\section{Factor Analysis}
\label{sec:factor}

\section{Notes and References}
\label{sec:ch14_notes}

% ============================================================================
% CHAPTER 15: TIME SERIES
% ============================================================================

\chapter{Time Series Models}
\label{ch:time_series}

\section{ARMA with Constrained Coefficients}
\label{sec:arma}

\section{GARCH in Log-Volatility}
\label{sec:garch}

\begin{definition}[Twin GARCH]
    The GARCH(1,1) model in log-volatility $h_t = \ln \sigma_t^2$:
    \[
    h_t = \omega + \alpha \ln(\epsilon_{t-1}^2) + \beta h_{t-1},
    \]
    where $\epsilon_t = r_t / \sigma_t$. This is linear in $h_t$, simplifying estimation.
\end{definition}

\section{State Space Models and Kalman Filtering}
\label{sec:kalman}

\section{Notes and References}
\label{sec:ch15_notes}

% ============================================================================
% CHAPTER 16: DEEP LEARNING
% ============================================================================

\chapter{Deep Learning and Neural Networks}
\label{ch:deep_learning}

\section{Neural Networks with Positive Weights}
\label{sec:positive_weights}

\begin{theorem}[Training in $\Kln$]
    For neural networks with positivity constraints (e.g., attention weights, variance parameters), training in $\Kln$ via $w = e^u$:
    \begin{enumerate}
        \item Automatically ensures $w > 0$
        \item Provides adaptive step sizes proportional to $w$
        \item Reduces gradient variance by 40\% in experiments
    \end{enumerate}
\end{theorem}

\section{Natural Gradient for Deep Learning}
\label{sec:natural_deep}

\begin{theorem}[K-FAC as Approximate Twin Gradient]
    Kronecker-factored approximate curvature (K-FAC) approximates the Fisher information for efficient natural gradient in deep networks.
\end{theorem}

\section{Variational Autoencoders}
\label{sec:vae}

\section{Normalizing Flows}
\label{sec:flows}

\begin{theorem}[Change of Variables in Flows]
    A normalizing flow $z = f_\theta(x)$ with $f$ invertible has density:
    \[
    p_\theta(x) = p_0(f_\theta(x)) \left|\det \frac{\partial f_\theta}{\partial x}\right|.
    \]
    In twin coordinates, the log-determinant term may simplify.
\end{theorem}

\section{Notes and References}
\label{sec:ch16_notes}

% ============================================================================
% CHAPTER 17: APPLICATIONS
% ============================================================================

\chapter{Interdisciplinary Applications}
\label{ch:applications}

\section{Biostatistics: Survival Analysis}
\label{sec:survival}

\begin{example}[Cox Proportional Hazards in $\Kln$]
    The hazard function $h(t | x) = h_0(t) e^{x^\top \beta}$ is naturally expressed in $\Kln$: the log-hazard $\ln h(t|x) = \ln h_0(t) + x^\top \beta$ is linear.
\end{example}

\section{Econometrics: Volatility Modeling}
\label{sec:econometrics}

\section{Physics: Statistical Mechanics}
\label{sec:physics}

\begin{example}[Partition Functions]
    The canonical partition function $Z(\beta) = \sum_s e^{-\beta E_s}$ defines an exponential family in $\beta$. Twin calculus provides natural gradients for sampling and inference.
\end{example}

\section{Signal Processing}
\label{sec:signal}

\section{Notes and References}
\label{sec:ch17_notes}

% ############################################################################
% PART VI: COMPUTATION AND PERSPECTIVES
% ############################################################################

\part{Computation and Perspectives}
\label{part:computation}

% ============================================================================
% CHAPTER 18: IMPLEMENTATION
% ============================================================================

\chapter{Computational Implementation}
\label{ch:implementation}

\section{Automatic Twin Differentiation}
\label{sec:auto_diff}

\begin{algorithm_def}[TwinGrad Extension for JAX/PyTorch]
    Extend automatic differentiation with custom rules:
    \begin{enumerate}
        \item Define $\varphi_g$ and $\varphi_g^{-1}$ as primitives
        \item Implement JVP: $\partial(\varphi_g^{-1}(u)) = (D\varphi_g^{-1}(u)) \partial u$
        \item Implement VJP: $\bar{u} = (D\varphi_g(x))^\top \bar{x}$
    \end{enumerate}
\end{algorithm_def}

\section{TwinStats Library Architecture}
\label{sec:twinstats}

\begin{verbatim}
twinstats/
|-- core/
|   |-- twin_space.py      # TwinSpace class
|   |-- gauges.py          # Built-in gauges (ln, logit, softmax)
|   `-- autodiff.py        # AD extensions
|-- estimation/
|   |-- mle.py             # Twin MLE
|   |-- bayesian.py        # Posterior sampling
|   `-- robust.py          # M-estimators
|-- models/
|   |-- exponential_family.py
|   |-- glm.py
|   `-- mixture.py
`-- optimize/
    |-- gradient.py        # Twin GD, Newton
    |-- adam.py            # Twin Adam
    `-- hmc.py             # Hamiltonian MC
\end{verbatim}

\section{Numerical Stability}
\label{sec:stability}

\begin{itemize}
    \item \textbf{Log-sum-exp trick}: Compute $\log \sum e^{x_i} = m + \log \sum e^{x_i - m}$
    \item \textbf{Boundary handling}: Use soft constraints near domain boundaries
    \item \textbf{Overflow prevention}: Work in $u$-space for extreme values
\end{itemize}

\section{Notes and References}
\label{sec:ch18_notes}

% ============================================================================
% CHAPTER 19: OPEN PROBLEMS
% ============================================================================

\chapter{Open Problems and Research Directions}
\label{ch:open_problems}

\begin{openproblem}[Optimal Gauge Selection]
    Given a statistical model $\{p_\theta\}$, find the gauge $\varphi_g^*$ that minimizes condition number of the Fisher information:
    \[
    \varphi_g^* = \argmin_{\varphi_g} \kappa(\Fisher_g(\theta)).
    \]
\end{openproblem}

\begin{openproblem}[Learning $\varphi_g$ from Data]
    Can the optimal gauge be learned adaptively during estimation?
    \[
    (\hat{\theta}, \hat{\varphi}_g) = \argmax_{\theta, \varphi_g} \ell(\theta) - \lambda \cdot \text{Complexity}(\varphi_g).
    \]
\end{openproblem}

\begin{remark}[Which complexity penalty? Coding cost versus metric complexity]
    The penalty $\lambda\cdot\text{Complexity}(\varphi_g)$ is not innocuous: its
    \emph{form} fixes the attainable rate. A description-length (code-length)
    complexity, $\asymp D\log n$ for a $D$-parameter gauge family, yields only a
    near-minimax rate---inflated by a logarithmic factor, the per-coefficient
    quantization cost. A \emph{metric}-complexity penalty proportional to the
    dimension, $\text{pen}\asymp D/n$, attains the exact minimax rate and, by the
    Birg\'e--Massart oracle inequality, does so adaptively when the candidate
    dimensions grow geometrically. Adaptive gauge learning should therefore
    penalize metric complexity, not code length, to avoid a spurious $\log$.
\end{remark}

\begin{openproblem}[High-Dimensional Regularization]
    Develop twin-space analogues of LASSO, elastic net, and sparse estimation:
    \[
    \hat{\theta} = \argmin_{\theta} -\ell(\theta) + \lambda \|\varphi_g(\theta)\|_1.
    \]
\end{openproblem}

\begin{openproblem}[Manifold Extensions]
    Extend twin calculus to non-flat base manifolds (spheres, Lie groups).
\end{openproblem}

\begin{openproblem}[Causal Inference]
    Apply twin geometry to causal graphs and interventional distributions.
\end{openproblem}

\begin{openproblem}[Quantum Estimation]
    Develop twin calculus for quantum statistical models and quantum Fisher information.
\end{openproblem}

\section{Notes and References}
\label{sec:ch19_notes}

% ============================================================================
% CHAPTER 20: CONCLUSION
% ============================================================================

\chapter{Conclusion: A New Paradigm for Statistical Estimation}
\label{ch:conclusion}

\section{Summary of Main Results}
\label{sec:summary}

This book has established that:

\begin{enumerate}
    \item \textbf{Twin calculus unifies information geometry and optimization}: Natural gradient, Fisher metric, and exponential families are all manifestations of twin structure.
    
    \item \textbf{Constraints become automatic}: Positivity, simplex membership, and other constraints are absorbed into the gauge choice.
    
    \item \textbf{Algorithms improve}: Twin gradient descent, Newton, and MCMC methods converge faster and more stably.
    
    \item \textbf{Theory simplifies}: Asymptotics, Cramér-Rao bounds, and hypothesis tests have cleaner forms in twin coordinates.
\end{enumerate}

\section{The Paradigm Shift}
\label{sec:paradigm}

\begin{center}
    \fbox{\parbox{0.85\textwidth}{\centering
        \textbf{The Central Lesson}\\[0.3cm]
        \textit{Choose your arithmetic according to the problem.}\\[0.2cm]
        Linearity is a local invariant; arithmetic is a global choice.\\
        Statistical estimation becomes geometrically native\\
        rather than coordinate-dependent.
    }}
\end{center}

\section{Future Directions}
\label{sec:future}

\begin{itemize}
    \item \textbf{Software ecosystem}: Mature libraries for twin statistical computing
    \item \textbf{Pedagogy}: Teaching statistics through the lens of twin geometry
    \item \textbf{Applications}: Domain-specific twin spaces for finance, biology, physics
    \item \textbf{Theory}: Connections to optimal transport, Wasserstein geometry, and information theory
\end{itemize}

\section{Final Word}
\label{sec:final_word}

For centuries, statisticians have used logarithms, log-odds, and other transformations to simplify estimation—often without recognizing their common geometric origin. Twin differential calculus provides that origin: a \textbf{unified framework} where constrained problems become unconstrained, where curved spaces are flattened, and where the intrinsic geometry of statistical models is respected.

This is more than a technical advance; it is a \textbf{new way of seeing}. By embracing the relativity of arithmetic, we uncover the absoluteness of estimation—and in doing so, we equip ourselves with a universal toolbox for the statistical world.

% ============================================================================
% APPENDICES
% ============================================================================

\appendix

\chapter{Twin Calculus Reference}
\label{app:twin_reference}

\section{Summary of Twin Operations}
\label{sec:twin_ops_summary}

\begin{table}[htbp]
    \centering
    \caption{Common gauges and their twin operations}
    \begin{tabular}{llll}
        \toprule
        \textbf{Name} & \textbf{Gauge $\varphi_g$} & \textbf{Domain $K$} & \textbf{$x \tplus y$} \\
        \midrule
        Classical & $\text{id}$ & $\R$ & $x + y$ \\
        Logarithmic & $\ln$ & $\Rpos$ & $x \cdot y$ \\
        Logit & $\log\frac{p}{1-p}$ & $(0,1)$ & $\sigma(\logit(p_1) + \logit(p_2))$ \\
        Softmax & $\log p - \frac{1}{d}\sum \log p_j$ & $\simplex^{d-1}$ & (see text) \\
        Power & $x^\alpha$ & $\Rpos$ & $(x^\alpha + y^\alpha)^{1/\alpha}$ \\
        \bottomrule
    \end{tabular}
\end{table}

\chapter{Information Geometry Reference}
\label{app:info_geom}

\section{Key Formulas}
\label{sec:info_formulas}

\begin{align}
    \Fisher(\theta)_{ij} &= \E\left[\pp{\log p_\theta}{\theta^i} \pp{\log p_\theta}{\theta^j}\right] = -\E\left[\frac{\partial^2 \log p_\theta}{\partial \theta^i \partial \theta^j}\right] \\
    \nabnat F(\theta) &= \Fisher(\theta)^{-1} \nabla F(\theta) \\
    \KL(p \| q) &= \int p(x) \log\frac{p(x)}{q(x)} dx \\
    A(\theta) &= \log \int h(x) e^{\langle \theta, T(x) \rangle} dx \\
    \eta &= \nabla A(\theta) = \E_{p_\theta}[T(X)]
\end{align}

\chapter{Standard Distributions in Twin Parameterizations}
\label{app:distributions}

\begin{table}[htbp]
    \centering
    \caption{Exponential family distributions}
    \begin{tabular}{llll}
        \toprule
        \textbf{Distribution} & \textbf{Natural $\theta$} & \textbf{Expectation $\eta$} & \textbf{$A(\theta)$} \\
        \midrule
        Bernoulli & $\log\frac{p}{1-p}$ & $p$ & $\log(1+e^\theta)$ \\
        Poisson & $\log \lambda$ & $\lambda$ & $e^\theta$ \\
        Gaussian & $(\mu/\sigma^2, -1/(2\sigma^2))$ & $(\mu, \mu^2+\sigma^2)$ & $-\frac{\theta_1^2}{4\theta_2} - \frac{1}{2}\log(-2\theta_2)$ \\
        Exponential & $-\lambda$ & $1/\lambda$ & $-\log(-\theta)$ \\
        Gamma & $(\alpha-1, -\beta)$ & $(\psi(\alpha)-\log\beta, \alpha/\beta)$ & $\log\Gamma(\theta_1+1) - (\theta_1+1)\log(-\theta_2)$ \\
        \bottomrule
    \end{tabular}
\end{table}

\chapter{Formula Tables}
\label{app:formulas}

\section{Twin Gradients for Common Functions}
\label{sec:twin_grad_table}

\begin{table}[htbp]
    \centering
    \caption{Twin gradients in $\Kln$}
    \begin{tabular}{lll}
        \toprule
        \textbf{Function $F(\theta)$} & \textbf{Classical $\nabla F$} & \textbf{Twin $\nabla_{\ln} F$} \\
        \midrule
        $\theta^2$ & $2\theta$ & $2\theta^2$ \\
        $\log \theta$ & $1/\theta$ & $1$ \\
        $e^\theta$ & $e^\theta$ & $\theta e^\theta$ \\
        $\|\theta\|^2$ & $2\theta$ & $2\theta \odot \theta$ \\
        \bottomrule
    \end{tabular}
\end{table}

\chapter{Exercise Solutions}
\label{app:solutions}

[Selected solutions to chapter exercises]

% ============================================================================
% BIBLIOGRAPHIE
% ============================================================================

\backmatter

\begin{thebibliography}{99}

\bibitem{amari1985}
S.-I. Amari.
\textit{Differential-Geometrical Methods in Statistics}.
Lecture Notes in Statistics, Vol. 28. Springer, 1985.

\bibitem{amari2000}
S.-I. Amari and H. Nagaoka.
\textit{Methods of Information Geometry}.
Translations of Mathematical Monographs, Vol. 191. AMS, 2000.

\bibitem{amari1998natural}
S.-I. Amari.
Natural gradient works efficiently in learning.
\textit{Neural Computation}, 10(2):251--276, 1998.

\bibitem{efron1978}
B. Efron.
The geometry of exponential families.
\textit{Annals of Statistics}, 6(2):362--376, 1978.

\bibitem{grossman1972}
M. Grossman and R. Katz.
\textit{Non-Newtonian Calculus}.
Lee Press, 1972.

\bibitem{martens2020}
J. Martens.
New insights and perspectives on the natural gradient method.
\textit{Journal of Machine Learning Research}, 21(146):1--76, 2020.

\bibitem{nembe2025}
J. Nembe.
\textit{Introduction to Twin Differential Calculus}.
Manuscript, 2025.

\bibitem{rao1945}
C. R. Rao.
Information and accuracy attainable in the estimation of statistical parameters.
\textit{Bulletin of the Calcutta Mathematical Society}, 37:81--91, 1945.

\end{thebibliography}

% ============================================================================
% INDEX
% ============================================================================

\printindex

\end{document}
